\documentclass[12pt,a4paper]{report}
\usepackage[utf8]{inputenc}
\usepackage[T1]{fontenc}
\usepackage[french]{babel}
\usepackage{graphicx}
\usepackage{hyperref}
\usepackage{booktabs}
\usepackage{geometry}
\usepackage{xcolor}
\geometry{margin=2.5cm}

\title{CyberGuard MLOps: plateforme de detection d'intrusions IoT avec CICIoT2023}
\author{Filiere CI 2 Informatique - IA \& BI-DS, Dev, Cyber}
\date{Mai 2026}

\begin{document}
\maketitle
\tableofcontents

\chapter{Introduction}
Ce rapport presente un projet MLOps complet pour la detection d'intrusions dans des flux reseau IoT. Le projet utilise le dataset reel CICIoT2023 et couvre l'ingestion, la validation, l'entrainement, le tracking MLflow, le deploiement FastAPI/Docker, le monitoring Prometheus/Grafana, la detection de drift et l'orchestration Airflow.

\chapter{Dataset CICIoT2023}
Le dataset CICIoT2023 a ete publie en 2023 par le Canadian Institute for Cybersecurity de l'University of New Brunswick. Il contient du trafic IoT issu de 105 appareils et 33 attaques regroupees en DDoS, DoS, Recon, Web-based, Brute Force, Spoofing et Mirai.

\chapter{Architecture MLOps}
\includegraphics[width=\textwidth]{../screenshots/v2_04_dvc_pipeline_ciciot2023.png}

\chapter{Validation, Tests et Qualite}
Le projet utilise Pytest, Black, isort, Ruff, MyPy, une suite Great Expectations documentee et des controles anti-fuite de donnees.

\chapter{Experimentation et Registry MLflow}
\includegraphics[width=\textwidth]{../screenshots/v2_06_mlflow_registry_ciciot2023.png}

\chapter{Deploiement et Monitoring}
\includegraphics[width=\textwidth]{../screenshots/v2_11_grafana_soc_dashboard.png}

\chapter{Resultats}
Les metriques finales sont extraites de \texttt{models/metrics.json}. Les figures finales doivent utiliser la nomenclature \texttt{screenshots/v2\_*.png} documentee dans \texttt{docs/SCREENSHOTS.md}.

\chapter{Conclusion}
CyberGuard demontre un cycle MLOps industrialise: donnees reelles, validation, experiment tracking, registry, serving, observabilite, drift monitoring et gouvernance.

\bibliographystyle{plain}
\bibliography{references}
\end{document}
